\chapter{Subprogram Calls\label{chap:SubprogramCalls}}

Subprograms can be called via an expression or as a statement.
Subprogram calls include invocations of getters and setters.

\ExampleDef{Recursive Subprogram Calls}
Subprograms can call each other recursively.
For example, in \listingref{MutualRecursion},
the longest possible call path is \verb|foo->bar->foo->bar->foo->bar->foo->bar|.
The final call to \verb|bar| is a \dynamicerrorterm.
\ASLListing{Mutual recursion}{MutualRecursion}{\definitiontests/MutualRecursion.bad.asl}

\ExampleDef{Invoking Getters and Setters}
\listingref{AccessorDeclaration} shows an example of declaring an accessor \verb|X|
and then invoking its getter and its setter.

\ChapterOutline
\begin{itemize}
  \item \FormalRelationsRef{Subprogram Calls} defines the formal relations for subprogram calls;
  \item \SyntaxRef{Subprogram Calls} defines the syntax of subprogram calls;
  \item \AbstractSyntaxRef{Subprogram Calls} defines the abstract syntax of subprogram calls;
  \item \TypeRulesRef{Subprogram Calls} defines the type rules for subprogram calls,
        and includes a discussion on parameter omission (\secref{ParameterOmission}); and
  \item \SemanticsRulesRef{Subprogram Calls} defines the dynamic semantics of subprogram calls.
\end{itemize}

%%%%%%%%%%%%%%%%%%%%%%%%%%%%%%%%%%%%%%%%%%%%%%%%%%%%%%%%%%%%%%%%%%%%%%%%%%%%%%%%
\FormalRelationsDef{Subprogram Calls}
\paragraph{Syntax:} Subprogram calls are grammatically derived from $\Nexpr$ and $\Nstmt$.
\paragraph{Abstract Syntax:} Subprogram calls are derived in the abstract syntax by
  $\ECall$ and $\SCall$;
\paragraph{Typing:} Subprogram calls are annotated by $\annotatecall$;
\paragraph{Semantics:} Subprogram calls are evaluated by $\evalcall$.

%%%%%%%%%%%%%%%%%%%%%%%%%%%%%%%%%%%%%%%%%%%%%%%%%%%%%%%%%%%%%%%%%%%%%%%%%%%%%%%%
\SyntaxDef{Subprogram Calls}
%%%%%%%%%%%%%%%%%%%%%%%%%%%%%%%%%%%%%%%%%%%%%%%%%%%%%%%%%%%%%%%%%%%%%%%%%%%%%%%%
\begin{flalign*}
\Nexpr \derives\  & \Tidentifier \parsesep \PlistZero{\Nexpr} &\\
\Nstmt \derives \ & \Tidentifier \parsesep \PlistZero{\Nexpr} \parsesep \Tsemicolon &
\end{flalign*}

%%%%%%%%%%%%%%%%%%%%%%%%%%%%%%%%%%%%%%%%%%%%%%%%%%%%%%%%%%%%%%%%%%%%%%%%%%%%%%%%
\AbstractSyntaxDef{Subprogram Calls}
%%%%%%%%%%%%%%%%%%%%%%%%%%%%%%%%%%%%%%%%%%%%%%%%%%%%%%%%%%%%%%%%%%%%%%%%%%%%%%%%
\RenderTypes[remove_hypertargets]{calls}


%%%%%%%%%%%%%%%%%%%%%%%%%%%%%%%%%%%%%%%%%%%%%%%%%%%%%%%%%%%%%%%%%%%%%%%%%%%%%%%%
\TypeRulesDef{Subprogram Calls}
%%%%%%%%%%%%%%%%%%%%%%%%%%%%%%%%%%%%%%%%%%%%%%%%%%%%%%%%%%%%%%%%%%%%%%%%%%%%%%%%

We define the following rules for annotating subprogram calls:
\begin{TypingRuleList}{SubprogramCalls}

\TypingRuleDef[annotates a subprogram call.]{AnnotateCall}
\RenderRelation{annotate_call}

\ExampleDef{Parameterised Subprogram Calls}
The specifications in \listingref{AnnotateCall} and \listingref{AnnotateCall2} are well-typed.
\ASLListing{Well-typed parameterised subprogram calls}{AnnotateCall}{\typingtests/TypingRule.AnnotateCall.asl}
\ASLListing{Well-typed parameterised subprogram calls}{AnnotateCall2}{\typingtests/TypingRule.AnnotateCall2.asl}

\ExampleDef{A Call with an Illegal Parameter Type}
The specification in \listingref{AnnotateCall-bad} is illegal since
the type of \verb|wid| --- \verb|integer{2,4,8}| --- does not \typesatisfyterm{}
the type of the parameter \verb|N| --- \verb|integer{2,4}|.
\ASLListing{An illegal parameter type}{AnnotateCall-bad}{\typingtests/TypingRule.AnnotateCall.bad.asl}

The specification in \listingref{AnnotateCall-bad2} is illegal since
expressions passed as parameters need to be \constrainedintegersterm.
\ASLListing{An illegal parameter type}{AnnotateCall-bad2}{\typingtests/TypingRule.AnnotateCall.bad2.asl}

The specification in \listingref{AnnotateCall-bad3} is illegal since
expressions passed as parameters need to be \constrainedintegersterm.
\ASLListing{An illegal parameter type}{AnnotateCall-bad3}{\typingtests/TypingRule.AnnotateCall.bad3.asl}

The specification in \listingref{AnnotateCall-bad4} is illegal since
the type of \verb|N| --- \verb|integer{N}| --- does not \typesatisfyterm{} \verb|integer{1,2,3}|.
\ASLListing{An illegal parameter type}{AnnotateCall-bad4}{\typingtests/TypingRule.AnnotateCall.bad4.asl}

For more examples, see \ExampleRef{Annotating Calls with Typed Parameters and Typed Arguments}.


\RenderProseAndFormally{annotate_call}
\CodeSubsection{\AnnotateCallBegin}{\AnnotateCallEnd}{../Typing.ml}


\TypingRuleDef[annotates a call with typed actual parameters and arguments.]{AnnotateCallActualsTyped}
\RenderRelation{annotate_call_actuals_typed}

\ExampleDef{Annotating Calls with Typed Parameters and Typed Arguments}
In \listingref{AnnotateCallActualsTyped-bad1}, the call expression \verb|xor_extend{64}(bv1, bv2)|
is ill-typed, since \\
\verb|xor_extend| has two parameters and only one was supplied.
\ASLListing{An ill-typed call expression}{AnnotateCallActualsTyped-bad1}{\typingtests/TypingRule.AnnotateCallActualsTyped.bad1.asl}

In \listingref{AnnotateCallActualsTyped-bad2}, the call expression \verb|xor_extend{64, 32}(bv1)|
is ill-typed, since \\
\verb|xor_extend| has two arguments and only one was supplied.
\ASLListing{An ill-typed call expression}{AnnotateCallActualsTyped-bad2}{\typingtests/TypingRule.AnnotateCallActualsTyped.bad2.asl}

In \listingref{AnnotateCallActualsTyped-bad3}, the call expression \verb|plus{64}(bv1, w)|
is ill-typed, since \\
the type of the argument \verb|w|, which is \verb|integer{0..128}|, does not \typesatisfyterm{}
the type of the formal argument \verb|z| (\verb|integer{0..N}|), with \verb|N| substituted by \verb|64|.
That is, \\
\verb|integer{0..64}|.
\ASLListing{An ill-typed call expression}{AnnotateCallActualsTyped-bad3}{\typingtests/TypingRule.AnnotateCallActualsTyped.bad3.asl}

In \listingref{AnnotateCallActualsTyped-bad4}, the call expression \verb|ones{myWid}()|
is ill-typed, since \verb|myWid| has the \unconstrainedintegertypeterm.
\ASLListing{An ill-typed call expression}{AnnotateCallActualsTyped-bad4}{\typingtests/TypingRule.AnnotateCallActualsTyped.bad4.asl}


\RenderProseAndFormally{annotate_call_actuals_typed}
\CodeSubsection{\AnnotateCallActualsTypedBegin}{\AnnotateCallActualsTypedEnd}{../Typing.ml}

\TypingRuleDef[inserts omitted parameters in calls to standard library subprograms.]{InsertStdlibParam}
\RenderRelation{insert_stdlib_param}

Note that this function relies on all standard library functions with input parameters having one of two simple forms:
\begin{lstlisting}
  func stdlibA{N}  (arg1: bits(N), ...) => ...
  func stdlibB{M,N}(arg1: bits(N), ...) => bits(...M...)
\end{lstlisting}

\ExampleDef{Inserting Parameters in Calls to Standard Library Subprograms}
The specification in \listingref{InsertStdlibParam}
shows examples of calls to standard library functions with some or all of their
parameters elided, and the equivalent calls with all parameters included.
\ASLListing{Inserting parameters in calls to standard library subprograms}{InsertStdlibParam}{\typingtests/TypingRule.InsertStdlibParam.asl}


\RenderProseAndFormally{insert_stdlib_param}

\TypingRuleDef[defines when a standard library parameter can be omitted.]{CanOmmitStdlibParam}
\RenderRelation{can_omit_stdlib_param}

\ExampleDef{Determining Whether a Parameter Can be Omitted}
The following are examples of subprogram signatures
where the first parameter can be omitted:
\begin{lstlisting}
func UInt{N: integer{1..128}} (x: bits(N)) => integer{0..2^N-1}
func Len{N}(x: bits(N)) => integer {N}
func ZeroExtend {N,M} (x: bits(M)) => bits(N)
\end{lstlisting}

The following are examples of subprogram signatures
where the first parameter cannot be omitted:
\begin{lstlisting}
func ReplicateBit{N}(isZero: boolean) => bits(N)
func Ones{N}() => bits(N)
\end{lstlisting}

\RenderProseAndFormally{can_omit_stdlib_param}

\TypingRuleDef[checks that actual parameter types satisfy formal parameter types.]{CheckParamsTypeSat}
\RenderRelation{check_params_typesat}

\ExampleDef{Checking that Expression Types Type-satisfy Parameters}
In \listingref{CheckParamsTypeSat},
annotating the call expression \verb|FlipSlice{FOUR, EIGHT}(bv)|
requires checking that both expression \verb|FOUR| and \verb|EIGHT| are \symbolicallyevaluableterm{}
and constrained, which they are.
Since the parameter \verb|M| is annotated with the type \\
\verb|integer{0..64}|,
the typechecker checks that the type of \verb|FOUR|, which is \verb|integer{4}| \typesatisfiesterm{}
\verb|integer{0..64}|.
The parameter \verb|N| is not annotated with a type --- its type is
determined to be \\
$\TInt(\Parameterized(\texttt{N}))$ --- and there is no need to check that
the type of \verb|EIGHT|, which is \verb|integer{8}|, \typesatisfiesterm{} $\TInt(\Parameterized(\texttt{N}))$.

\ASLListing{Checking that expression types type-satisfy parameters}{CheckParamsTypeSat}{\typingtests/TypingRule.CheckParamsTypeSat.asl}


\RenderProseAndFormally{check_params_typesat}

\TypingRuleDef[substitutes actual parameter expressions into parameterised types.]{RenameTyEqs}
\RenderRelation{rename_ty_eqs}

\ExampleDef{Transforming Parameterised Types Based on Actual Arguments}
In \listingref{RenameTyEqs}, annotating the call expression \verb|FlipPrefix{8}(bv, 4)|
requires substituting \verb|8| for \verb|N| in the types
\verb|bits(N)| and \verb|integer{0..N}|, yielding the types
\verb|bits(8)| and \verb|integer{0..8}|.
\ASLListing{Transforming parameterised types based on actual arguments}{RenameTyEqs}{\typingtests/TypingRule.RenameTyEqs.asl}


\RenderProseAndFormally{rename_ty_eqs}

\TypingRuleDef[substitutes and normalizes parameter expressions.]{SubstExprNormalize}
\RenderRelation{subst_expr_normalize}

\ExampleDef{Substituting Parameter Expressions and Normalising}
In \listingref{SubstExpr},
considering the call expression \verb|plus{z + 22}(bv1, z)|,
the expression \verb|z + 22| is substituted for the parameter \verb|N| in the type
\verb|bits(N + 2)| used for the argument \verb|x| and as the return type of \verb|plus|,
resulting in the expression \verb|(40 + 22) + 2|, which is then normalised into \verb|64|.
\ASLListing{Substituting Parameter Expressions}{SubstExpr}{\typingtests/TypingRule.SubstExpr.asl}


\RenderProseAndFormally{subst_expr_normalize}

\TypingRuleDef[substitutes parameter expressions in expressions.]{SubstExpr}
\RenderRelation{subst_expr}

The function assumes that $\ve$ appears in the declaration of a parameter,
which means it is in the subset allowed by $\extractparameters$.

See \ExampleRef{Substituting Parameter Expressions and Normalising}.


\RenderProseAndFormally{subst_expr}
\CodeSubsection{\SubstExprBegin}{\SubstExprEnd}{../ASTUtils.ml}

\TypingRuleDef[substitutes parameter expressions in constraints.]{SubstConstraint}
\RenderRelation{subst_constraint}

\ExampleDef{Substituting Parameter Expressions in Constraints}
In \listingref{SubstExpr},
considering the call expression \verb|plus{z + 22}(bv1, z)|,
the expression \verb|z + 22| is substituted for the parameter \verb|N| in the type
\verb|integer{0..N}| used for the argument \verb|z|,
resulting in the type \verb|integer{0..64}|.


\RenderProseAndFormally{subst_constraint}

\TypingRuleDef[checks that actual argument types satisfy formal argument types.]{CheckArgsTypeSat}
\RenderRelation{check_args_typesat}

\ExampleDef{Checking that Actual Arguments Type-satisfy the Formal Arguments}
In \listingref{RenameTyEqs}, checking that the call \verb|FlipPrefix{8}(bv, 4)| is well-typed
requires checking that the types of the arguments --- \verb|bits(8)| and \verb|integer{4}| ---
\typesatisfyterm{} the corresponding types of the formal arguments --- \verb|bits(N)| and \verb|integer{0..N}| ---
once the parameter \verb|N| is substituted for \verb|8|.
That is, that \verb|bits(8)| \typesatisfiesterm{} \verb|bits(8)| and that
\verb|integer{4}| \typesatisfiesterm{} \verb|integer{0..8}|.
Since both these checks hold, the call is indeed well-typed.

In contrast, a call \verb|FlipPrefix{8}(bv, 9)| is ill-typed since
\verb|integer{9}| does not \typesatisfyterm{} \verb|integer{0..8}|.

We note that in the following, \TypingRuleRef{AnnotateCallActualsTyped} guarantees
that $\funcsigargs$ and $\argtypes$ have the same length.

\RenderProseAndFormally{check_args_typesat}

\TypingRuleDef[annotates the return type of a subprogram call.]{AnnotateRetTy}
\RenderRelation{annotate_ret_ty}

\ExampleDef{Annotating the Return Type of a Subprogram Call}
In \listingref{AnnotateRetTy}, annotating the return type \verb|bits(N)| of the subprogram \verb|flip|
for call expression \verb|flip{64}(bv)| yields the annotated type \verb|bits(64)|.

Since \verb|proc| does not have a return type, the call statement \verb|proc()| does not require
annotating a return type (thus $\annotateretty$ yields $\none$ for $\rettyopt$).
\ASLListing{Annotating the Return Type of a Subprogram Call}{AnnotateRetTy}{\typingtests/TypingRule.AnnotateRetTy.asl}

\listingref{AnnotateRetTy-bad} shows an example of an ill-typed call statement \verb|flip{64}(bv);|
and an ill-typed call expression \verb|proc()|, since procedures can only be used in call statements
and functions can only be used in call expressions.
\ASLListing{Ill-typed Subprogram Calls}{AnnotateRetTy-bad}{\typingtests/TypingRule.AnnotateRetTy.bad.asl}


\RenderProseAndFormally{annotate_ret_ty}

\TypingRuleDef[identifies the subprogram matching a call signature.]{SubprogramForSignature}
\RenderRelation{subprogram_for_signature}

This function determines which one of
the following cases holds:
\begin{enumerate}
\item there is no declared subprogram that matches $\name$, $\callerargtypes$, \\ and $\calltype$; or
\item there is exactly one subprogram that matches $\name$, $\callerargtypes$, \\ and $\calltype$;
\end{enumerate}
If more than one subprogram that matches $\name$ and
$\callerargtypes$, this is detected by the rule
\TypingRuleRef{DeclareSubprograms}, which invokes the
rule \\ \TypingRuleRef{DeclareOneFunc}, which invokes
the rule \TypingRuleRef{AddNewFunc}, which results in
a \typingerrorterm{}.

The first case results in a \typingerrorterm{}.
If the second case holds, the function returns a tuple which comprises:
\begin{itemize}
\item $\namep$ --- the string that uniquely identifies this subprogram;
\item $\callee$ --- the AST node defining the called subprogram; and
\item $\vses$ --- the set of \sideeffectdescriptorsterm{} associated with $\name$.
\end{itemize}
\ProseOtherwiseTypeError

\ExampleDef{Matching a Subprogram to an Identifier}

\listingref{SubprogramForSignature} shows an example where all subprogram calls
match subprogram declarations.
\ASLListing{Successfully matching subprogram signatures to definitions}{SubprogramForSignature}{\typingtests/TypingRule.SubprogramForSignature.asl}

\listingref{SubprogramForSignature-undefined} shows an example where
the subprogram call \verb|add_10(5)| is illegal, since no subprogram named \verb|add_10|
is declared.
\ASLListing{Subprogram signature undefined}{SubprogramForSignature-undefined}{\typingtests/TypingRule.SubprogramForSignature.bad.undefined.asl}

\listingref{SubprogramForSignature} shows an example where
the subprogram call \verb|add_10(5.0)| is illegal, since, although a subprogram named \verb|add_10|
is declared, it does not match the required signature (the type of the first argument is
the \integertypeterm{} rather than the \realtypeterm).
\ASLListing{Subprogram name does not match signature}{SubprogramForSignature-no-candidates}{\typingtests/TypingRule.SubprogramForSignature.bad.no_candidates.asl}


\RenderProseAndFormally{subprogram_for_signature}
\CodeSubsection{\SubprogramForSignatureBegin}{\SubprogramForSignatureEnd}{../Typing.ml}

\TypingRuleDef[filters subprogram candidates for a call signature.]{FilterCallCandidates}
\RenderRelation{filter_call_candidates}

The names $\candidates$ are assumed to exist in $\tenv.\staticenvsG.\subprograms$.

\ExampleDef{Filtering Subprograms Matching a Call}
In \listingref{SubprogramForSignature}, filtering the
set of declared subprograms for the call expression\\ \verb|add_10(5)| yields
the (singleton list containing the) subprogram \\
\verb|func add_10(x: integer) => integer|.
In contrast, filtering the set of declared subprograms for the call expression
\verb|add_10(5.0)| in \listingref{SubprogramForSignature-no-candidates},
yields an empty list.


\RenderProseAndFormally{filter_call_candidates}

\TypingRuleDef[defines compatible subprogram and call types.]{CallTypeMatches}
\RenderRelation{call_type_matches}

\ExampleDef{Matching call types}
The table below shows keywords used to define functions, and their possible compatible call types.

\begin{center}
\begin{tabular}{ll}
\textbf{Keyword used to define function} & \textbf{Compatible call types} \\
\hline
\verb|func|                            & $\STFunction$, $\STProcedure$ \\
\verb|getter| (inside \verb|accessor|) & $\STFunction$, $\STGetter$    \\
\verb|setter| (inside \verb|accessor|) & $\STSetter$                   \\
\end{tabular}
\end{center}


\RenderProseAndFormally{call_type_matches}

\CodeSubsection{\CallTypeMatchesBegin}{\CallTypeMatchesEnd}{../Typing.ml}

\TypingRuleDef[checks whether an argument type clashes with a list of argument types.]{HasArgClash}
\RenderRelation{has_arg_clash}

\ExampleDef{Argument Clashing}
\listingref{type-clashes} shows examples of types ---
the arguments of procedures -- that do not clash and types that do clash
(shown in comments).
\ASLListing{Examples of Argument clashing}{type-clashes}{\typingtests/TypingRule.TypeClashes.asl}


\RenderProseAndFormally{has_arg_clash}
\CodeSubsection{\HasArgClashBegin}{\HasArgClashEnd}{../Typing.ml}

\TypingRuleDef[defines type clashing for subprogram arguments.]{TypeClash}
\RenderRelation{type_clashes}


Note that \Prosetypeclashing{} is an equivalence relation.
In particular note that if $T$ \Prosetypeclashes{} with $A$ and $B$ then $A$ and
$B$ \Prosetypeclash{}.

See \ExampleRef{Argument Clashing}

\ExampleDef{Ill-typed Subprogram Declarations}
In specification \listingref{TypeClash-bad}, the type \verb|integer| \emph{\Prosetypeclashes} \verb|time|
in the \staticenvironmentterm{} where both types have been annotated,
which is why both declarations of \verb|f| clash and thus they are ill-typed.
\ASLListing{Ill-typed Subprogram Declarations}{TypeClash-bad}{\typingtests/TypingRule.TypeClashes.bad.asl}


\RenderProseAndFormally{type_clashes}
\CodeSubsection{\TypeClashBegin}{\TypeClashEnd}{../types.ml}

\subsubsection{Comment}
Note that if $\vt$ \subtypesatisfiesterm{} $\vs$ then $\vt$ and $\vs$ \Prosetypeclash, but not the other
way around.

\TypingRuleDef[annotates a list of expressions.]{ExpressionList}
\RenderRelation{annotate_exprs}

\ExampleDef{Annotating a List of Expressions}
In \listingref{ExpressionList}, the list of expressions \verb|10, 20.0|
is annotated as
\[
\left[
\begin{array}{l}
\left(\begin{array}{l}\TInt(\WellConstrained([\AbbrevConstraintExact{\ELInt{10}}])),\\ \ELInt{10},\\ \emptyset\end{array}\right),\\
\left(\begin{array}{l}\TReal,\\ \eliteral{\LReal(20)},\\ \emptyset\end{array}\right)
\end{array}
\right]
\]
\ASLListing{Annotating a list of expressions}{ExpressionList}{\typingtests/TypingRule.ExpressionList.asl}


\RenderProseAndFormally{annotate_exprs}

\end{TypingRuleList}

\subsection{Parameter Omission\label{sec:ParameterOmission}}

ASL allows dropping a parameter from the parameter list of a subprogram call
during type checking of standard library calls (\TypingRuleRef{InsertStdlibParam}).
In this case, a parameter expression can be inferred by matching the call expression
against one of two patterns (see the rule for details and examples).
If omitting a parameter yields an empty list, the \verb|{...}| must be entirely removed.
We refer to this as \emph{parameter omission}.

This situation is exemplified in \listingref{ParameterOmission}.
\ASLListing{Parameter omission}{ParameterOmission}{\definitiontests/ParameterOmission.asl}

The specification in \listingref{ParameterOmission-bad} is ill-typed,
as an empty parameter list for a standard library function must be entirely omitted.
\ASLListing{Erroneous parameter omission}{ParameterOmission-bad}{\definitiontests/ParameterOmission.bad.asl}

%%%%%%%%%%%%%%%%%%%%%%%%%%%%%%%%%%%%%%%%%%%%%%%%%%%%%%%%%%%%%%%%%%%%%%%%%%%%%%%%
\SemanticsRulesDef{Subprogram Calls}
%%%%%%%%%%%%%%%%%%%%%%%%%%%%%%%%%%%%%%%%%%%%%%%%%%%%%%%%%%%%%%%%%%%%%%%%%%%%%%%%

We define the following rules for evaluating subprogram calls:
\begin{SemanticsRuleList}{SubprogramCalls}

\SemanticsRuleDef[evaluates a subprogram call.]{Call}
\RenderRelation{eval_call}

The evaluation first evaluates the expressions corresponding to the arguments
and parameters and then passes their values in a resulting configuration
to the helper relation $\evalsubprogram$.

\ExampleDef{Calling Subprograms}
In \listingref{Call}, calling \verb|non_throwing_func| terminates normally,
while calling \\
\verb|throwing_func| terminates by throwing an exception.
\ASLListing{Calling subprograms}{Call}{\semanticstests/SemanticsRule.Call.asl}


\RenderProseAndFormally{eval_call}
\CodeSubsection{\EvalCallBegin}{\EvalCallEnd}{../Interpreter.ml}

\SemanticsRuleDef[evaluates a subprogram body for a call.]{EvalSubprogram}
\RenderRelation{eval_subprogram}

\ExampleDef{Subprogram Calls}
In \listingref{semantics-fcall},
the function \texttt{main} calls the function \texttt{foo} and the procedure \texttt{bar}.
\ASLListing{Evaluating subprogram calls}{semantics-fcall}{\semanticstests/SemanticsRule.FCall.asl}


\RenderProseAndFormally{eval_subprogram}
\CodeSubsection{\EvalFCallBegin}{\EvalFCallEnd}{../Interpreter.ml}

\subsubsection{Comments}

It is not an error for execution of a procedure or setter to end without a
return statement.

\SemanticsRuleDef[checks the recursion limit for a subprogram call.]{CheckRecurseLimit}
\RenderRelation{check_recurse_limit}

\ExampleDef{Checking the Recursion Limit of a Subprogram Call}
In \listingref{CheckRecurseLimit-nolimit}, the function \verb|factorial| is specified without a limit expression,
and evaluating \verb|main| terminates normally.
\ASLListing{A recursive function with no limit expression}{CheckRecurseLimit-nolimit}{\semanticstests/SemanticsRule.CheckRecurseLimit.no_limit.asl}

In \listingref{CheckRecurseLimit}, the function \verb|factorial| specifies the limit expression \verb|11|,
and evaluating \verb|main| terminates normally.
\ASLListing{A recursive function with a limit expression}{CheckRecurseLimit}{\semanticstests/SemanticsRule.CheckRecurseLimit.asl}

In \listingref{CheckRecurseLimit-limitreached}, the function \verb|factorial| specifies the limit expression \verb|10|,
and evaluating \verb|main| terminates with a \dynamicerrorterm{} (\LimitExceeded), since the limit is exceeded.
\ASLListing{A recursive call exceeding the specified limit}{CheckRecurseLimit-limitreached}{\semanticstests/SemanticsRule.CheckRecurseLimit.limit_reached.asl}


\RenderProseAndFormally{check_recurse_limit}

\SemanticsRuleDef[reads the value produced by a subprogram call.]{ReadValueFrom}
\RenderRelation{read_value_from}

% NO_EXAMPLE


\RenderProseAndFormally{read_value_from}
\CodeSubsection{\EvalReadValueFromBegin}{\EvalReadValueFromEnd}{../Interpreter.ml}

\SemanticsRuleDef[assigns actual argument values to formal arguments.]{AssignArgs}
\RenderRelation{assign_args}

\ExampleDef{Assigning Arguments}
In \listingref{AssignArgs}, the call expression \verb|plus(10, 5)| binds \verb|x| to $\nvint(10)$
and $\verb|y|$ to $\nvint(5)$. The expression \verb|x + y| then accesses
these variable values and evaluates to $\nvint(15)$.
\ASLListing{Assigning arguments}{AssignArgs}{\semanticstests/SemanticsRule.AssignArgs.asl}


\RenderProseAndFormally{assign_args}

\SemanticsRuleDef[converts subprogram-return configurations into call results.]{MatchFuncRes}
Using the type \RenderType{TContinuingOrReturning} we define the following relation.

\RenderRelation{match_func_res}
\ExampleDef{Converting Configurations Upon Subprogram Return}
In \listingref{MatchFuncRes},
the final configuration resulting from evaluating the call \verb|proc()|
is a \Prosecontinuingconfiguration{}, which is converted into
a \Prosenormalconfiguration{} with no return value.
On the other hand, the configuration resulting from evaluating the
call \verb|returns_values()| is a \Prosereturningconfiguration{} with the value
$\nvint(5)$, which is converted into a \Prosenormalconfiguration{}
with the same value and a fresh identifier for it.

\ASLListing{Converting configurations upon subprogram return}{MatchFuncRes}{\semanticstests/SemanticsRule.MatchFuncRes.asl}


\RenderProseAndFormally{match_func_res}

\end{SemanticsRuleList}
